\documentclass[10pt]{article}
\usepackage[T1]{fontenc}
\usepackage[utf8]{inputenc}
\usepackage[french]{babel}
\usepackage{amsmath,amssymb,mathrsfs,mathpazo}
\usepackage{geometry}
\geometry{a4paper,margin=2cm}
\usepackage{graphicx,array,tabularx,tabularray,colortbl,multirow,multicol,booktabs,enumitem,pifont,ulem,fancybox,diagbox,cancel}
\usepackage{pstricks,pstricks-add,pst-plot,pst-tree,pst-node,pst-text,pst-func,pst-math,pst-bspline,pst-3dplot,pst-eucl,pst-solides3d}
\let\pscalculate\undefined
\usepackage{pst-bezier}
\usepackage{tikz,pgfplots,tkz-tab}
\usetikzlibrary{arrows,arrows.meta,patterns,calc,positioning,trees,decorations.markings,shapes.geometric,shapes.misc,shapes.arrows,intersections,decorations.pathmorphing,decorations.pathreplacing,backgrounds,shadows,babel}
\pgfplotsset{compat=1.18}
\usepgfplotslibrary{fillbetween}
\usepackage[np]{numprint}
\usepackage{listings,makecell,pas-tableur,twoopt}
\let\barre\undefined
\usepackage[pstricks]{tablvar}
\definecolor{darkgreen}{RGB}{0,100,0}
\providecommand{\red}{\color{red}}
\providecommand{\blue}{\color{blue}}
\providecommand{\green}{\color{green}}
\providecommand{\black}{\color{black}}
\providecommand{\white}{\color{white}}
\providecommand{\gray}{\color{gray}}
\providecommand{\cyan}{\color{cyan}}
\providecommand{\magenta}{\color{magenta}}
\providecommand{\yellow}{\color{yellow}}
\providecommand{\violet}{\color{violet}}
\providecommand{\orange}{\color{orange}}
\providecommand{\euro}{\texteuro}
\providecommand{\EUR}{\texteuro}
\providecommand{\ieme}{\textsuperscript{e}}
\providecommand{\degres}{\ensuremath{^\circ}}
\providecommand{\degree}{\ensuremath{^\circ}}
\providecommand{\vv}[1]{\overrightarrow{#1}}
\providecommand{\Coord}[1]{#1}
\providecommand{\decosix}{}
\providecommand{\starredbullet}{\ensuremath{\bullet}}
\pagestyle{empty}
\setlength{\parindent}{0pt}
\usepackage{adjustbox}\DeclareUnicodeCharacter{2236}{:}\DeclareUnicodeCharacter{2190}{\ensuremath{\leftarrow}}
\providecommand{\R}{\mathbb{R}}
\providecommand{\N}{\mathbb{N}}
\providecommand{\Z}{\mathbb{Z}}
\providecommand{\Q}{\mathbb{Q}}
\providecommand{\D}{\mathbb{D}}
\providecommand{\C}{\mathbb{C}}
\providecommand{\e}{{\mathrm{e}}}
\providecommand{\vect}[1]{\overrightarrow{#1}}
\providecommand{\vectt}[1]{\overrightarrow{\mathrm{#1}}}
\providecommand{\vvt}[1]{\overrightarrow{\mathrm{#1}}}
\providecommand{\np}[2][]{{#2~#1}}
\providecommand{\vv}[1]{{\overrightarrow{#1}}}
\renewcommand{\d}{{\mathrm{d}}}
\providecommand{\degres}{{^\circ}}
\providecommand{\degre}{{^\circ}}
\providecommand{\blue}{\color{blue}}
\providecommand{\red}{\color{red}}
\providecommand{\enskip}{\quad}
\providecommand{\llbracket}{\lbrack\!\lbrack}
\providecommand{\rrbracket}{\rbrack\!\rbrack}
\providecommand{\euro}{{\text{€}}}
\providecommand{\no}{{\text{n°}}}
\providecommand{\interval}[2]{{[#1;#2]}}
\providecommand{\intervFO}[2]{{[#1;#2[}}
\providecommand{\dotfill}{\ldots}

\providecommand{\ds}{\displaystyle}
\providecommand{\textding}[1]{#1}
\providecommand{\Oij}{\(\mathrm{O};\vec{\imath},\vec{\jmath}\)}
\providecommand{\Oijk}{\(\mathrm{O};\vec{\imath},\vec{\jmath},\vec{k}\)}
\providecommand{\Ouv}{\(\mathrm{O};\vec{u},\vec{v}\)}
\providecommand{\intervFF}[2]{\ensuremath{[#1;#2]}}
\renewcommand{\euro}{\ifmmode\text{\texteuro}\else\texteuro\fi}
\DeclareUnicodeCharacter{00B0}{\ifmmode^{\circ}\else\textdegree\fi}
\DeclareUnicodeCharacter{20AC}{\ifmmode\text{\texteuro}\else\texteuro\fi}
% Énoncés : Éducation nationale. Transcription et corrigés : APMEP, Denis Vergès et contributeurs. Découpage et classement : CoopMaths, Sébastien Lozano, Éric Elter et contributeurs. Adaptation de la présentation : Maths Explorer.
% CC BY-SA : https://coopmaths.fr/www/about/
% https://coopmaths.fr/alea/static/bac/2023/tex/bac_2023_03_sujet1_asie_2_cor.tex
\begin{document}
\noindent\textit{APMEP / CoopMaths. CC BY-SA. Présentation adaptée par Maths Explorer.}\par\medskip
\medskip

\begin{enumerate}
\item On a les coordonnées suivantes :
\quad
	$\vect{\text{MN}}\begin{pmatrix} x_\mathrm{N} - x_\mathrm{M}\\ y_\mathrm{N}-y_\mathrm{M} \\ z_\mathrm{N} - z_\mathrm{M} \end{pmatrix}=\begin{pmatrix} -1\\ -\frac12 \\ \frac14 \end{pmatrix}$ et $\vect{\text{MP}}\begin{pmatrix}0\\-1\\-2\end{pmatrix}$.

\item On obtient la figure suivante :
	
	
	
	
	
	
	
	
	
	
	
	

\hfill~\begin{tikzpicture}[x={(50mm,-7mm)},y={(32mm,11mm)},z={(0mm,57mm)}]
		\coordinate (A) at (0,0,0); \coordinate (E) at (0,0,1);
		\coordinate (B) at (1,0,0); \coordinate (F) at (1,0,1);
		\coordinate (C) at (1,1,0); \coordinate (G) at (1,1,1);
		\coordinate (D) at (0,1,0); \coordinate (H) at (0,1,1);
		\coordinate (M) at (1,1,0.75); \coordinate (N) at (0,0.5,1);
		\coordinate (P) at (1,0,-1.25); 

		\draw (A)node[below left]{A}--(B)node[below]{B}
			--(C) node[right] {C} -- (G) node[above right]{G}
			--(H) node[above] {H} -- (E) node[left]{E}--cycle
			(B) -- (F) node[above]{F}--(E) (F)--(G)  ;
		\draw [dash pattern= on 2mm off 1mm]
		(A) -- (D) node[below]{D}--(H) (D)--(C)  ;
		\draw (M)++(-2pt,-2pt)--++(4pt,4pt) ++(-4pt,0pt)--++(4pt,-4pt) node[right]{M};
		\draw (N)++(-2pt,-2pt)--++(4pt,4pt) ++(0,-4pt)--++(-4pt,4pt) node[above]{N};
		\draw (P)++(-2pt,-2pt)--++(4pt,4pt) ++(0,-4pt)--++(-4pt,4pt) node[left]{P};
		\draw [dotted](B)--(P);
	\end{tikzpicture}\hfill~

\item Les vecteurs $\vect{\mathrm{MN}}$ et $\vect{\mathrm{MP}}$ sont clairement non colinéaires (l'abscisse de l'un est nulle sans que celle l'autre ne le soit), donc les trois points M, N et P sont non alignés, et définissent donc un plan.

\item
	\begin{enumerate}
		\item Puisque le repère est orthonormé, on va calculer le produit scalaire avec les coordonnées :

$\vect{\text{MN}} \cdot \vect{\text{MP}} = (-1)\times 0 + \left(\dfrac{-1}{2}\right) \times (-1) + \dfrac{1}{4}\times (-2) = 0 + \dfrac{1}{2} - \dfrac{1}{2} = 0 $.

Les vecteurs sont donc orthogonaux et non nuls, donc l'angle qu'ils définissent est un angle droit. On en déduit que le triangle est rectangle en M.

Calculons les longueurs MN et MP :

$\mathrm{MN} = \sqrt{\left(-1\right)^2+\left(-\dfrac{1}{2}\right)^2+\left(\dfrac{1}{4}\right)^2 } = \dfrac{\sqrt{21}}{4}$.
\qquad
$\mathrm{MP} = \sqrt{\left(0\right)^2+\left(-1\right)^2+\left(-2\right)^2} = \sqrt{5}$

Comme $\mathrm{MN} \neq \mathrm{MP}$, le triangle est seulement un triangle rectangle en M (il n'est pas isocèle).

	\item L'aire du triangle MNP est la moitié du produit de la longueur d'une base par la longueur de la hauteur correspondante.

Comme le triangle est rectangle en M, on choisit [MN] comme base et donc [MP] est donc la hauteur correspondante.

L'aire du triangle est donc $\mathcal{A}_{\mathrm{MNP}} = \dfrac{\dfrac{\sqrt{21}}{4}\times \sqrt{5}}{2} = \dfrac{\sqrt{105}}{8}$
	\end{enumerate}
\item
	\begin{enumerate}
		\item Avec $\vect{n}(5~;~-8~;~4)$ on a :

$\vect{\mathrm{MN}}\cdot \vect{n} = (-1)\times 5 + \left(-\dfrac{1}{2}\right) \times (-8) + \dfrac{1}{4} \times 4 = -5 + 4 + 1 = 0$

$\vect{\mathrm{MP}}\cdot \vect{n} = 0\times 5 + (-1) \times (-8) + (-2) \times 4 = 0 + 8 -8 = 0$

Le vecteur non nul $\vect{n}$ est donc orthogonal à deux vecteurs non colinéaires du plan (MNP), il est donc orthogonal au plan.

C'est donc un vecteur normal de (MNP).

\item On en déduit que le plan (MNP) a une équation cartésienne de la forme $5x - 8y + 4z + d= 0$, où $d$ est un nombre réel.

Comme $\mathrm{M} \in (\mathrm{MNP})$, $d$ est tel que :

$\begin{aligned} 5x_\mathrm{M} - 8y_\mathrm{M} + 4z_\mathrm{M} + d= 0
&\iff 5\times 1 - 8\times 1 + 4\times \dfrac{3}{4} + d =0\\
&\iff 5 - 8 + 3 + d=0\\
&\iff d = 0\\ \end{aligned}$

Une équation cartésienne de (MNP) est donc bien $5x-8y+4z=0$.
	\end{enumerate}
\item Si une droite est orthogonale au plan (MNP), c'est donc qu'elle est dirigée par un vecteur normal à (MNP), donc $\vect{n}$, par exemple.

Le point F a pour coordonnées F(1~;~0~;~1) et $\vect{n}$ a pour coordonnées $\vect{n}(5~;~-8~;~4)$ et donc la droite passant par F et dirigée par $\vect{n}$ admet pour représentation paramétrique le système : $\left\{\begin{array}{l} x=x_{\mathrm{F}} + tx_{\vect{n}}\\
	y=y_{\mathrm{F}} + ty_{\vect{n}}\\
	z=z_{\mathrm{F}} + tz_{\vect{n}} \end{array}\right.\quad t \in \R $,\qquad  soit, ici : $\left\{\begin{array}{l} x=1 + 5t\\
	y=\phantom{1}-8t\\
	z=1+4t \end{array}\right.\quad t \in \R $

\item Si L le projeté orthogonal du point F sur le plan (MNP), alors, c'est que L est l'intersection du plan (MNP) avec la droite $d$, orthogonale à (MNP) et passant par F.

Si on considère le point $L_t$, de paramètre $t$ sur la droite $d$, cherchons la valeur du paramètre $t$ pour que $L_t \in (\mathrm{MNP})$.

$ \begin{aligned} L_t \in (\mathrm{MNP}) &\iff 5x_{L_t} - 8y_{L_t} + 4z_{L_t} = 0\\
		&\iff 5(1 + 5t) - 8(-8t) + 4(1 + 4t) = 0\\
		&\iff 5 + 25t +64t + 4 + 16 t = 0\\
		&\iff 105t + 9 = 0\\
		&\iff t = \dfrac{-9}{105} = \dfrac{-3}{35}	\\ \end{aligned}$

Le point de paramètre $\dfrac{-3}{35}$ sur $d$ a pour coordonnées : $\left(1 + 5\times \dfrac{-3}{35}~;~-8\times \dfrac{-3}{35}~;~1 + 4\times \dfrac{-3}{35}\right) $, c'est donc bien le point L, de coordonnées  L$\left(\dfrac47~~;~~\dfrac{24}{35}~~;~~\dfrac{23}{35}\right)$.

\item On a donc $\vect{\mathrm{FL}}\begin{pmatrix} \frac{-3}7\\ \frac{24}{35}\\ \frac{-12}{35}	\end{pmatrix}$ et donc FL $= \sqrt{\left(\dfrac{-3}{7}\right)^2 + \left(\dfrac{24}{35}\right)^2 + \left(\dfrac{-12}{35}\right)^2 } = \dfrac{3\sqrt{105}}{35}$.

Pour calculer le volume du tétraèdre FMNP, on va considérer le triangle MNP comme étant la base, avec $\mathcal{A}_{\mathrm{MNP}} =\dfrac{\sqrt{105}}{8} $.

La hauteur correspondante est donc la distance du quatrième sommet F au plan contenant la base, c'est-à-dire le plan (MNP). D'après la question \textbf{7.}, c'est donc la distance FL, que l'on vient de calculer.

On a donc $ V_{\mathrm{FMNP}} = \dfrac{1}{3} \times \mathcal{A}_{\mathrm{MNP}} \times \mathrm{FL} = \dfrac{1}{3}\times \dfrac{\sqrt{105}}{8} \times \dfrac{3\sqrt{105}}{35} = \dfrac{105}{8\times 35} = \dfrac{3\times 35}{8\times 35 } = \dfrac{3}{8}$

Le volume de FMNP est donc de $\dfrac{3}{8}$.
\end{enumerate}

\medskip

\end{document}
